% !TeX program = xelatex
% 在 docs/latex/ai-authoring 中用 XeLaTeX 编译 npc-companion.tex。
\documentclass[UTF8,11pt,fontset=fandol,a4paper]{ctexart}
\usepackage[a4paper,top=22mm,bottom=22mm,left=23mm,right=23mm]{geometry}
\usepackage{tikz,tabularx,booktabs,array,fancyhdr,hyperref}
\usetikzlibrary{arrows.meta,fit,backgrounds}
\definecolor{ink}{HTML}{183540}\definecolor{teal}{HTML}{126B62}
\definecolor{muted}{HTML}{465F69}\definecolor{pale}{HTML}{E9F4F0}
\hypersetup{colorlinks=true,urlcolor=teal,linkcolor=teal,pdftitle={Redflag 工坊搭档：NPC 玩法与引擎计划}}
\ctexset{section={format=\Large\bfseries\color{ink},afterindent=true}}
\pagestyle{fancy}\fancyhf{}\lhead{\small Redflag\enspace /\enspace 工坊搭档}\rhead{\small 玩法与引擎设计\enspace ·\enspace 2026.11.22}\rfoot{\small\thepage}
\setlength{\parindent}{2\ccwd}\setlength{\parskip}{.54\baselineskip}\setlength{\emergencystretch}{4em}
\newcommand{\file}[1]{\texttt{\detokenize{#1}}}
\newcommand{\src}[2]{\href{#1}{#2}}
\newcommand{\page}[1]{\clearpage\section{#1}}
\newcommand{\captext}[1]{\par{\small\color{muted}\noindent #1\par}}
\renewcommand{\arraystretch}{1.35}
\newcolumntype{L}[1]{>{\raggedright\arraybackslash}p{#1}}
\newcolumntype{Y}{>{\raggedright\arraybackslash}X}
\begin{document}
\section{跟一个记得你的技师开一次工坊}
玩家和技师阿榆要在开门前挂好一块旧招牌。玩家只能清楚看见招牌正面，扶着背面的阿榆能感觉到自己是否还托着它。谁去找工具、谁扶哪一边、是否保留旧牌上的划痕，都由两人商量。收尾时门打开，两个人一起退一步，看自己挂好的招牌。第二局换一间新扫描的工坊，上一局磨合出来的递工具习惯仍然有用。这是面向一至两名真人和一个 NPC 的终版玩法设想。

\textbf{第一分钟，玩家自然学会和它说话。} 玩家喊「帮我拿胶带缠过的那把锤子」。阿榆抬头循着玩家的指向看，两把都像，木柄都在盒沿后面。它移到桌角看见一蓝一黑，问「蓝的这把？」同时用手指向那把。玩家说「嗯，给我」，它拿起并把手柄朝外递到玩家够得着的位置。玩家不说话，也可以用指向、点头或接过来继续。若仅有一个清楚符合条件的目标，就省下这句追问。

\textbf{第二件事让玩家看见性格。} 阿榆喜欢修旧东西。它看见旧牌上的划痕，趁放下工具的空隙问「这道留着？以前也是这样挂的吧。」台词只引用眼前划痕；关于以前的事是一个问句。玩家说「留着，这是我爸刻的」。以后这块牌再搬家，它拿布时会绕过这个位置。若玩家答「磨掉吧」，最终画面就保留经过打磨的版本。阿榆可以有偏好，最后的决定属于玩家。

\textbf{第三件事让搭档关系有用。} 玩家双手托着牌说「等一下」，阿榆停止递锤，留在已承托的那一侧。若玩家改口「我扶着，你看歪没歪」，它先等玩家接稳，再走到正面看。玩家可以自己装完，叫阿榆只在要倾倒时提醒。结束后，它把灯放到玩家实际约定的位置。游戏用招牌的姿态、旧划痕和工具归位处表现两人的选择。

\textbf{正式内容的规模。} 建议围绕同一工坊先设计三种可复玩的委托：挂旧招牌、根据一张局部照片帮访客认领旧物、停电后合力把物件摆回原位。三者分别检验身体协作、多轮辨认和记忆更新。新场景可以换道具、材质、遮挡与动线；每局至少保留一个玩家能改变的结果。认领委托里，玩家决定是先对照片的缺口，还是先让访客描述旧物；停电委托中，提灯的人能看清标签，搬箱的人必须给同伴报位置。三个委托换的是双方各自掌握的信息与可以做的事。上述人物和玩法均为拟议内容。

\page{阿榆怎样决定看、问、动手和等待}
每次决策只解决眼下的一件事。听到「找蓝胶带那把」后，系统建立这次请求，记下说话人、那个时刻的指向和它自己看到的桌面。视觉模型找出两处像锤子的区域，交给 NPC 两个临时引用。它比较哪条线索能区分两者。可以走一步看见手柄，就先走；两侧仍都符合，就问玩家能区分它们的问题；已有明确证据且伸手可达，就提出拿取。相机观察到单个物件也可能看错，特别是颜色受反光影响时，要允许收回判断。

\begin{center}
\begin{tikzpicture}[x=1cm,y=1cm,n/.style={rounded corners=2pt,draw=teal!60,fill=white,text width=39mm,align=center,minimum height=11mm,font=\small},ar/.style={-{Stealth[length=2mm]},teal,thick}]
\path[draw=ink!50,rounded corners=3pt] (-.13,.78) rectangle (15.27,-4.15);
\fill[teal!7] (0,.64) rectangle (15.12,-1.47);
\fill[blue!5] (0,-1.61) rectangle (15.12,-4.01);
\node[anchor=west,font=\small\bfseries,text=ink] at (.12,.38){NPC 当下能知道什么};
\node[anchor=west,font=\small\bfseries,text=ink] at (.12,-1.88){游戏怎样兑现或拒绝这个动作};
\node[n](obs) at (2.38,-.64){眼睛、声音与指向\\带时刻和说话人};
\node[n](belief) at (7.54,-.64){眼前线索与共同约定\\记录来源及上次所见};
\node[n](choice) at (12.68,-.64){选择下一步\\看、问、帮忙或等待};
\node[n](next) at (2.38,-3.15){更新计划和表现\\玩家看见停手或递交};
\node[n](world) at (7.54,-3.15){权威游戏世界\\确认持有者和动作结果};
\node[n](act) at (12.68,-3.15){执行当前一步\\导航、伸手或说话};
\draw[ar](obs)--(belief);\draw[ar](belief)--(choice);\draw[ar](choice)--(act);\draw[ar](act)--(world);\draw[ar](world)--(next);
\end{tikzpicture}
\end{center}
\captext{图 1\quad 原创拟议调用图。箭头表示上一步把观察、动作请求或结果交给下一步。世界只把 NPC 可以感知的变化和它自己动作的回执交还给它，其他房间的真相只给服务器和测试裁判。}

\textbf{谁有权做哪件事。} LLM 解释指代、从角色当前证据提出下一步，并组织简短台词。C++ 负责可走、可见、能否伸手，以及持有与交接。角色运动、视线和手部每帧继续更新。只在话轮结束、目标失效、发现新线索或动作有结果时重新决策。技能用同一约定返回进行中、完成、中止或具体失败原因；重复的拿取请求只结算一次。外出找东西的计划可以挂起，回来时先确认位置和承诺是否仍有效。\src{https://say-can.github.io/}{SayCan，语言建议与技能可行性}

\textbf{遇到冲突时。} 玩家要求停下，立即取消排队动作；手里承重的物件留在稳定状态。玩家只插一句闲话时，阿榆可以停下台词而继续托住招牌。两名玩家同时要它帮不同的忙，阿榆记住已经答应谁，不把新请求悄悄覆盖旧请求。它可说「我先扶着，你来接一下？」第二人接稳的服务器回执到达后，才松手并继续。玩家能据此预测它，也能纠正它。

\page{主动帮忙的分寸与可玩的剧情}
阿榆的主动性主要表现在怎么做事。它看到玩家双手已占，先把任务需要的台灯转向工作处。要接走玩家手里的旧牌，先伸手示意，等双方交接。玩家两次绕开推荐的摆法，它就让路。玩家明确说「我来」后，本任务的同类建议立即静音；喊它名字或新的承托请求可以重新叫回它。若玩家主动要安静，就保留动作提示。设计目的在于让熟人知道它会怎么帮，初次来的人也能轻松拒绝。

\begin{center}\small
\begin{tabularx}{\linewidth}{@{}L{31mm}YY@{}}\toprule
发生了什么 & 玩家看见什么 & 游戏根据什么决定下一步\\\midrule
线索互相矛盾 & 它放低准备抓取的手，侧身看手柄；实在看不清才问 & 另一合法视角能否区分候选、玩家指示是否仍指向当前物件。\\
玩家把「蓝短柄」叫“小蓝” & 当局稍后玩家说「把小蓝放我左边」，它照做 & 把称呼与玩家、确认过的对象和这间工坊关联；场景变化后重新辨认实体。\\
玩家忙着对齐招牌 & 它扶住并保持安静；若要松手，先等人接稳 & 承托动作回执、玩家拒绝提示的记录、招牌是否已有第二个支点。\\
阿榆拿错了 & 玩家说「另一把」，它放回原处，指向新候选核对 & 刚才被排除的候选和修正均写入本轮记录；不重复同一个问题。\\
收工还有时间 & 它指着那道保留下来的划痕，邀玩家一起把旧牌摆正 & 只在有共同经历且玩家空闲时触发；拒绝后继续收工。\\\bottomrule
\end{tabularx}
\end{center}

\textbf{把记忆写成可用的约定。} 「刚才蓝柄在桌角」来自视觉，失去视线后只能当线索。「把工具放我左边」附说话人以及本次或以后都适用的范围。「我已经把锤子给你」要有真实交接回执。当前任务记候选、失败尝试和挂起步骤；存档只保留玩家明确教的称呼、偏好及共同完成的事件。给新玩家帮忙时不套用老玩家的偏好。玩家说「以后改右边」就覆盖旧约定；手工读档时人物记忆随该存档恢复。\src{https://arxiv.org/abs/2304.03442}{Generative Agents，记忆与规划的实验}

\textbf{写情境的人决定什么。} 每个剧情片段保存触发条件、玩家可选择的回应、可见后果、退出条件与重复间隔。例如「旧牌划痕」要等两人真的看见划痕、当前无人承重；两条选择都能完成修理。若玩家已经把牌挂好，就跳过这次闲聊。作者可把新片段写入场景项目；LLM 在当前事实范围里配话，导演只按共同经历与闲忙挑片段。这样同一关能留下不同回忆。\src{https://cdn.aaai.org/ojs/18722/18722-52-22361-1-10-20210928.pdf}{Mateas 与 Stern，Façade 的故事片段与联合行为}\quad\src{https://steamcdn-a.akamaihd.net/apps/valve/2009/ai_systems_of_l4d_mike_booth.pdf}{Valve 演讲中的同伴与戏剧节奏}

\page{从它眼前的一把锤子生成可移动物件}
普通模型关与扫描关使用同一套 NPC 指向、确认和抓取规则。模型关保留引擎已有的对象边界，但向 NPC 隐藏名字与层级。扫描关只给未命名的整屋 3DGS。离线可准备高斯外观、相机标定、导航和简化碰撞；按你的选项，第一次看到锤子时才开始找物件边界。能从同一个摄像机画面讨论桌面，是两条路径都有的能力。

\textbf{从「那把」到同一个目标。} 玩家指向时记下玩家相机、输入时刻和指向落点；NPC 侧只得知它当时实际能见到或玩家向它共享的指示。若只看到玩家的手势，先把它变成大致方向，走到能看到那张桌子的位置。视觉给出临时图像区域；第二个视角优先选择能分清手柄的可达位置。玩家给一张照片时可以匹配外观，但照片本身不能证明目标现在的房间位置。视角有遮挡就重看，若两把确实一样，问用途或让玩家指向某一把。\src{https://sony.github.io/B3-Seg-project/}{B³-Seg，按预期信息增益选取观察点}

\textbf{从「看见」到实际递出。} 从掩码和相机位姿计算高斯成员，剔除属于桌面的污染；用另一视角检查轮廓和接触边。服务端核对源资产版本，创建实例，装配简化碰撞、可达抓点，把选中部分从原屋外观和碰撞中移出，再广播成员与变换。递出时玩家接过的同一帧更换持有者。编辑器可以显示成员、抓点和碰撞，允许人修正或另存为可复用资产。源照片从未见过的背面或被盖桌面标为未知，可用另一次空桌采集作对照。重建算法的补画单独标记。\src{https://github.com/07LEE/gs-object-extraction}{高斯物件提取与独立 PLY 导出}

\textbf{在线划分超时时玩家仍有事可做。} 先指给玩家看，证据足够就报告「就在这儿」，需要对方拿起时才启动完整提取。超出等待预算时保留线索，允许玩家继续挂招牌、换物件或者让 NPC 稍后再来。提取结束后重新核对场景与任务版本。已被玩家取消的任务可以保留提取结果供编辑器查看，不能触发一次迟到的抓取。GaussianCut 已示范从单图提示得到高斯前景，作者指出图割有时要数分钟。\src{https://arxiv.org/html/2411.07555v1}{GaussianCut，第 6 节与局限}

\textbf{怎样衡量选型本身。} 另做一份同场景合并网格，给它同样的视觉模型、合法视角及在线分割机会。用未参与调参的工坊比较准备人时、首个可用物件的等待、误带桌面比例、放回后的残影、GPU 帧耗时与显存。纯在线先跑。若超预算，定位是绘制、视觉推理还是三维分割；只有后两项确实需要时才分别试预计算视觉特征和未命名的候选分组。每次只改变一项。网格与高斯可以独立赢得适合自己的内容范围。

\page{从 Delivery 接回 Redflag 的游戏循环}
\begin{center}\small
\begin{tabularx}{\linewidth}{@{}L{27mm}YY@{}}\toprule
责任 & 当前能对照的真实入口 & Redflag 拟建的接口及其输出\\\midrule
眼睛与观察 & \file{SceneRenderer::RenderRuntime} 收相机投影与变换；\file{Scene::OnUpdateRuntime} 选游戏主相机 & 为技师建立按需观察相机，返回带时刻、机位和场景版本的图；每个机位持有自己的历史缓冲。\\
决定与反馈 & Delivery 的 \file{BuildDecisionRequest} 汇集依据，\file{TryAcceptProposal} 核对请求和修订 & NPC 提出注视、移动、问、持有或交接的一步；世界返回成功或具体失败。日志留选择依据和引用的观察。\\
声音和身体 & Delivery 的玩家打断测试区分旧回复和独立活动；Redflag 场景启动当前接 2D 物理与脚本 & 轮次、动作任务、持有分别给编号。动作中止先问是否已承重。导航、骨骼姿势、三维碰撞与抓点接入后分别验证。\\
制作和试玩 & \file{ForgeLayer::OnScenePlay} 复制编辑场景，停 Play 回到编辑场景 & 人和制作助手共用可撤销编辑。试玩发现的部件由人选择「另存为资产」，场景里可选、可看、可改。\\\bottomrule
\end{tabularx}
\end{center}

\textbf{语音回路。} 游戏客户端经 WebRTC 将每名玩家的声音交给语音服务。STT 的临时结果可带动转头和听的姿势，动作请求使用已确认的轮次和说话人。LLM 按需取得 NPC 的观察与相关承诺，输出结构化建议，以及当下可说的话。TTS 先播一段有用的回应，直到玩家打断或轮次结束。玩家插话可立即压低或停止旧播报；单纯“嗯”与“停，先给我”要分开处理。停止播报、取消排队动作和释放手中物件分别有不同的触发条件。已确认的世界动作继续履约。LiveKit Agents 可用于编排，首批用语音和文本两种入口跑同一套委托。\src{https://docs.livekit.io/agents/logic/turns/}{LiveKit 的话轮与打断}

\textbf{并发与网络。} 推理在工作进程异步运行，游戏线程维护本帧世界和动画。请求携带轮次、任务、观察版本及说话人，回复过期就重看。推理离线时，能暂停的任务进入等待，有承托任务则继续承重，玩家仍能用指向和接物动作。两客户端争抢物件由服务器裁决；可靠确认谁持有，连续位姿用带序号的快照。相邻的 REDFServer 已核对到 TCP 接收入口，UDP 传输与游戏复制都列为待建。\src{https://github.com/ValveSoftware/GameNetworkingSockets}{Valve GameNetworkingSockets}

\captext{源码线索。Redflag 的场景启动见 \file{Scene.cpp} 第 135 行，运行相机见 \file{SceneRenderer.h} 第 88 行，试玩见 \file{ForgeLayer.cpp} 第 1226 行。Delivery 的请求与校验见 \file{DeliveryNpcAgentComponent.cpp} 第 615 与 2081 行，语音打断见 \file{DeliveryNpcPlayerInterruptionTests.cpp} 第 53 行，服务端流式入口见 \file{Tools/DeliveryAI/server.py} 第 428 行。}

\page{作者怎样制作一个会改变结果的场景}
作者输入一句话，例如「在这个工坊里做一局挂旧招牌的双人委托。搭档能扶、能找玩家指出的工具，玩家可以决定是否保留划痕。」制作助手向正在运行的编辑器查询可用动作、场景修订、资产、脚本和试玩入口。查到什么能用，再安排相机、出生点、支点、可行走区域、角色骨架和通关规则。剧情可以定义「招牌挂好并且玩家开门」这种世界条件。测试裁判知道答案，但给 NPC 的输入只含合法观察，尤其不能把材质名、层级名或目标索引塞进它的上下文。

\textbf{把人物和动作交给不同的作者。} 人物设定保存名字、旧物偏好、何时主动、被拒绝后怎样退开，以及用哪些身体动作表达。技师的动作库支持注视、指向、行走、让路、等待、托住、拿、递、放和简短发言。人可以在属性面板调整递物高度、抓点、拒绝提示的作用范围等；两位作者分别编辑条件和台词时，按字段合并。让技师学会一句新暗号属于游戏内容。要让手指稳定停在世界某一点，则要增加或改进骨骼、IK 或动画控制，再重编引擎、重跑夹具。Python 可做离线动画重定向和运动匹配数据库准备，C++ 做逐帧求姿，游戏 C\# 只管委托规则；语言为建议分工。

\textbf{AI 接错该在哪里修。} 已有编辑器动作能表达的东西，用工具写入并回读。完成新委托若只缺一个获胜条件，改游戏脚本。若需要创建 NPC 独立相机或支持双人承托而引擎缺对应能力，就给出触发用例、提出局部 C++ 变更，在独立构建中验证后再回到委托。人在 AI 试玩时挪桌、换素材或手改脚本，助手对编辑器内最新修订和磁盘上的实际源码重新核对，只重算失效步骤。玩法提示与验证条件可写成开源 Skill 常用的步骤清单，MCP 传递引擎支持的结构化操作。Unity MCP 的公开 Skill 可作工具选择和回读方式的参照。\src{https://github.com/CoplayDev/unity-mcp}{Unity MCP 与作者提供的 Skill}

\textbf{人能调试一次没接稳的招牌。} 在工坊编辑器里选中那次失败，时间线并排显示玩家说「我接着」的时刻、技师亲眼看见的画面、当前承诺、请求、服务端判定、手部姿势与最终支点。点任一事件能跳回相应画面或调用点。评测记录结构化结果及可公开的行为理由，例如「两把均为蓝柄，选择询问长度」。调试只需要输入、选择和实际结果。设备、模型与提示版本、场景版本、总用时和 GPU 帧成本要随录像保存，否则换模型后难以复查。TEACh 的双视角回放与独立评测可作参考。\src{https://github.com/alexa/teach}{TEACh 作者仓库与回放入口}

\page{工程推进与每个阶段的过关条件}
以下为同一终版的依赖顺序。单人开发宜先串起可回放的主线；扫描分割可从阶段二起并行做实验。每段按具体验收结束，再估下一段工期。角色素材、目标 GPU、语音部署方式及联机并发均未定，当前没有可靠的总周数。

\begin{center}\small
\begin{tabularx}{\linewidth}{@{}L{24mm}YY@{}}\toprule
交付阶段 & 完成的工作 & 可以进入下一阶段的证据\\\midrule
1．可玩的对照关 & 用独立网格搭工坊、旧牌和两把相似工具；定好任务、玩家接手、重开与输入回放。先用固定策略扮演技师。 & 同一局用两种分工都能开门。能重放一次完成和一次失败，退出试玩不改编辑场景。\\
2．有身体的搭档 & 上线合法观察、别名、视角指代、一步动作选择、可达导航与基本骨骼；先用文本验规则，语音随后走同一入口。 & 两把相似时能侧看或问。朝向交换仍理解「我左边」；录像中玩家能看出它指哪把。\\
3．协作与剧情 & 加承诺、协同承重、递物回执、偏好存档、话轮打断；完成三个委托的触发、分歧和不同结尾。 & 托住时能被打断台词并保持支撑。更改偏好、拒绝建议、读档重来都可复现。\\
4．陌生工坊 & 对未分组高斯做在线发现；成员与三维碰撞可复查；另存为资产后再次载入。用合并网格作第二对照。 & 在保留工坊完成发现、分出、递交、另存；记录错误与人时，裁决哪些扫描关值得采用。\\
5．双人一致性 & 接语音服务、丢包延迟模拟、可靠状态与不可靠位姿、重连恢复。 & 两人争一件物品始终只有一个持有者。走开、重连、过期响应和中断均能正确收尾。\\
6．一句话出关卡 & 接编辑器能力查询、Skill、游戏脚本和独立构建，跑真人试玩、性能回归与发布包验收。 & 换一间房、换一条新委托，人可手改，AI 能接续。交付可启动游戏、可编辑项目和回放。\\\bottomrule
\end{tabularx}
\end{center}

\textbf{先后顺序可以调整的部分。} 语音可在动作夹具完成后尽早接入，交替听两人抢话的实际效果。运动匹配用同一段接近和让路的回放验证脚滑、停步和 CPU 开销；递物的最后一段还需抓点、姿态修正及真实交接帧。高斯提取跑得慢时，网格工坊仍可测角色与剧情。盲测证明提取对制作或玩法有收益后，再安排静态高斯在编辑器与运行时的长期维护。

\page{玩法验收：它能做对，玩家也想再玩}
\textbf{先留能判对错的回放。} 测八种情况：清楚且唯一的工具；玩家与 NPC 相向说「我左边」；同色工具需要一句追问；视线外的工具被挪走；别名被纠正；抓前被另一人拿走；承托时打断台词；玩家拒绝后又重新招呼。每种各置三种布局，用固定输入和录制的模型回复重复验证执行器；换真实模型时每种布局独立运行五次，保留全部失败与分母。测试与调参使用不同房间。物件暂时不可见且没有人提供线索时，正确结果可以是再观察、合理询问或说明暂时找不到。\src{https://aclanthology.org/2024.emnlp-main.1100/}{Tang 等关于三维环境中不同观察视角的指代研究}

\textbf{建议的首轮验收线。} 确定性回放要求旧话轮不能触发抓取、交接不产生双持、承托不能凭一句话丢失支点，违规即修复。真实模型先看可解题的最终正确交接率是否达到 90\%，分布局列出结果。建议目标机上「确认打断至旧语音停止」P95 在 250 毫秒内，「话轮结束至首条与任务有关的可听回应」P95 在 2 秒内；另报起始语音到检测打断的时间，避免只计已检测时段。把首次观察、三维提取、抓取提交、第二个客户端可见各自计时。空泛的“正在处理”不计有效回应。这些数值是拟议的调试门槛，需用首轮真人等待反馈校准，本项目尚无实测。

\textbf{好玩请人来判。} 先找 8 至 12 名玩家做小规模交叉试玩。一组先玩「只依请求办事」的搭档，一组先玩加入主动观察、记忆与性格选择的搭档。保持外观、声线、动作能力和语音等待条件一致，轮流分配两个难度相近的工坊。每局后先让玩家复述「哪一刻觉得它懂你」「哪次帮倒忙」，再看愿不愿意自选它玩额外一局。记录玩家主动教昵称、主动交托任务、打断建议和要求它离开的时刻。逐人保留选择与原因。这个样本用于发现问题与形成下一次实验假设，随后逐项关闭记忆、主动帮忙或剧情片段，找出差别来自哪里。自拟问项用于本项目反馈。玩家动机研究对自主感、能力感与联结的区分，可帮助归类访谈。\src{https://selfdeterminationtheory.org/wp-content/uploads/2020/10/2006_RyanRigbyPrzybylski_MandE.pdf}{Ryan 等，玩家动机研究}

\textbf{决定修哪个系统。} 能准确找到但总打断人，先调主动介入和话轮。每局都问同一个问题，查记忆是否记录了纠正、图像是否真的分得清。网格关好评而扫描关差，先看等候和抓错是否集中在分割阶段。若成功率高但没人愿意再玩，就改委托的分工与后果，让玩家能指出「因为我说了什么，它后来做了什么」。人觉得角色好玩、回放证明角色靠谱、扫描素材节省制作工时，这三项各自独立记录。

\captext{交付状态：本文为终版设计与验收计划。参考项目、论文与本地入口按正文核对。工坊玩法、NPC 推理、三维交互和各项指标未在 Redflag 中实施或实测。}
\end{document}
